% Options for packages loaded elsewhere
\PassOptionsToPackage{unicode}{hyperref}
\PassOptionsToPackage{hyphens}{url}
\documentclass[
]{article}
\usepackage{xcolor}
\usepackage{amsmath,amssymb}
\setcounter{secnumdepth}{-\maxdimen} % remove section numbering
\usepackage{iftex}
\ifPDFTeX
  \usepackage[T1]{fontenc}
  \usepackage[utf8]{inputenc}
  \usepackage{textcomp} % provide euro and other symbols
\else % if luatex or xetex
  \usepackage{unicode-math} % this also loads fontspec
  \defaultfontfeatures{Scale=MatchLowercase}
  \defaultfontfeatures[\rmfamily]{Ligatures=TeX,Scale=1}
\fi
\usepackage{lmodern}
\ifPDFTeX\else
  % xetex/luatex font selection
\fi
% Use upquote if available, for straight quotes in verbatim environments
\IfFileExists{upquote.sty}{\usepackage{upquote}}{}
\IfFileExists{microtype.sty}{% use microtype if available
  \usepackage[]{microtype}
  \UseMicrotypeSet[protrusion]{basicmath} % disable protrusion for tt fonts
}{}
\makeatletter
\@ifundefined{KOMAClassName}{% if non-KOMA class
  \IfFileExists{parskip.sty}{%
    \usepackage{parskip}
  }{% else
    \setlength{\parindent}{0pt}
    \setlength{\parskip}{6pt plus 2pt minus 1pt}}
}{% if KOMA class
  \KOMAoptions{parskip=half}}
\makeatother
\setlength{\emergencystretch}{3em} % prevent overfull lines
\providecommand{\tightlist}{%
  \setlength{\itemsep}{0pt}\setlength{\parskip}{0pt}}
\usepackage{bookmark}
\IfFileExists{xurl.sty}{\usepackage{xurl}}{} % add URL line breaks if available
\urlstyle{same}
% fallback for those not using the hyperref driver hyperxmp:
\makeatletter
\@ifundefined{xmpquote}{\newcommand{\xmpquote}[1]{#1}}{}
\makeatother
\hypersetup{
  hidelinks,
  pdfcreator={LaTeX via pandoc}}

\author{}
\date{}

\begin{document}

\section{Informe de puerta E6 --- run\_v1.22\_r2 (Paper IV
v1.22-r2)}\label{informe-de-puerta-e6--run_v122_r2-paper-iv-v122-r2}

Veredicto actual: \textbf{PASS\_WITH\_OBSERVATIONS}. Evidencia:
\texttt{20\_EVIDENCE/E6/}. Auditor: Claude Opus 5.5, misma sesión que
las ejecuciones anteriores; no se ejecutó Lean.

\begin{center}\rule{0.5\linewidth}{0.5pt}\end{center}

\section{E6 --- Correcciones españolas, paridad y PDF
(run\_v1.22\_r2)}\label{e6--correcciones-espauxf1olas-paridad-y-pdf-run_v122_r2}

Objetivo: delta ES de \texttt{v1.22\_editorial\_candidate} (r1) a
\texttt{v1.22\_editorial\_candidate\_r2} (r2). El inglés se trata en E0
(byte-idéntico, no regenerado).

\subsection{Método (independiente de los scripts del
editor)}\label{muxe9todo-independiente-de-los-scripts-del-editor}

\begin{enumerate}
\def\labelenumi{\arabic{enumi}.}
\tightlist
\item
  \textbf{Diff propio.}
  \texttt{20\_EVIDENCE/E6/auditor\_diff\_es\_md.diff} y
  \texttt{auditor\_diff\_es\_tex.diff}, más un inventario palabra a
  palabra (\texttt{10\_SCRIPTS/e6\_r2\_delta.py} →
  \texttt{E6\_R2\_DELTA.json}).

  \begin{itemize}
  \tightlist
  \item
    MD: 2565 → 2565 líneas, con 11 líneas y 13 cambios de palabra.
  \item
    TeX: 2367 → 2367 líneas, con 11 líneas y 13 cambios, en las mismas
    posiciones lógicas.
  \item
    El \texttt{CHANGES\_es.diff} del editor es idéntico byte a byte al
    diff unificado que regeneré (13717 caracteres).
  \end{itemize}
\item
  \textbf{Contenido protegido}, comparado entre r1 y r2:

  \begin{itemize}
  \tightlist
  \item
    Idénticos: 237 fórmulas en display, 1989 fórmulas en línea, 5
    bloques de código, 146 etiquetas \texttt{\textbackslash{}tag} y la
    sección de referencias completa.
  \item
    Códigos en línea: pasan de 330 a 331; ninguno se elimina y el único
    que aumenta es \texttt{\textasciigrave{}Model\textasciigrave{}}.
  \item
    La comparación de «líneas de encabezado o entrada en negrita» da
    \emph{false} por una sola razón: la línea 578 empieza con la entrada
    en negrita «Segunda fase\ldots», cuyo texto en negrita no cambia; sí
    cambia la palabra «matchings» del resto del párrafo.
  \item
    \texttt{IMPACT\_CHECK.json}: ninguna línea cambiada está en un
    párrafo de enunciado (Teorema, Lema, Proposición o Corolario), en
    una tabla ni en la bibliografía. Una cae dentro de una demostración
    (l.688, Lema 5.3, «ni otro empaquetamiento»).
  \end{itemize}
\item
  \textbf{El PDF corresponde al TeX corregido.}

  \begin{itemize}
  \tightlist
  \item
    El PDF r2 (sha256 8ebfccd7\ldots, 73 páginas, producido por
    xdvipdfmx a partir de LaTeX el 2026-10-01 10:50 UTC) se comparó con
    el r1 por texto extraído en tokens.
  \item
    El cuerpo de r1, aplicando exactamente las sustituciones del
    inventario, coincide con el cuerpo de r2 salvo en 3 diferencias.
    Todas son tipográficas y las verifiqué visualmente:

    \begin{itemize}
    \tightlist
    \item
      1 y 2, p. 55: la línea «La absorción da mín(u\_w,v\_w) ≤
      \textbar M\_w\textbar+t\_W+L\ldots» de r1 tenía el espaciado
      matemático comprimido; en r2 se reparte normalmente. La fórmula es
      la misma.
    \item
      3, p. 65--66: en r1, la Tabla 9 se partía entre páginas y repetía
      el encabezado «Enunciado / Declaración»; en r2 cabe en p. 65, así
      que el encabezado aparece una sola vez. Las 10 filas están
      presentes.
    \end{itemize}
  \item
    La bibliografía del PDF es idéntica.
  \item
    El log de compilación no tiene Overfull, «Missing character»,
    referencias indefinidas ni errores. Hay 5 avisos de Underfull, ya
    declarados por el editor.
  \item
    El preámbulo TeX (40 líneas) es idéntico al de r1.
  \end{itemize}
\item
  \textbf{Raster.} Comparando r1 y r2 a 72 dpi cambian 18 páginas: 16,
  18, 20, 21, 35, 55, 56, 61--71. Es exactamente la lista de
  \texttt{ARTIFACT\_CHECKS.json} y de
  \texttt{VISUAL\_REVIEW\_v1.22\_r2.md}. Las demás 55 páginas son
  raster-idénticas. El texto cambia en 17 páginas: la 35 cambia solo en
  la fuente de \texttt{Model}.
\item
  \textbf{Inspección visual} (\texttt{PAGE\_INSPECTION\_LOG.csv}, 73/73
  páginas):

  \begin{itemize}
  \tightlist
  \item
    Las 73 páginas r2 se revisaron en hojas 2-up a 90 dpi
    (\texttt{pairs/}).
  \item
    Las páginas 55, 56 y 61--71 se compararon además lado a lado
    r1\textbar r2 a 110 dpi (\texttt{cmp/}), mirando los
    desplazamientos.
  \item
    Las páginas 16, 18, 20, 21 y 35 se verificaron en las hojas r2. Sus
    imágenes \texttt{cmp/} no llegaron a mostrarse (C3-02). La p. 35 se
    vio además a 140 dpi (\texttt{zoom/es\_r2\_p035.png}).
  \item
    Regla de registro: cada fila se escribió solo con la imagen visible
    en ese momento. Las filas descartadas se conservan en
    \texttt{PAGE\_INSPECTION\_LOG.discarded\_unseen\_cmp.csv}.
  \end{itemize}
\end{enumerate}

\subsection{X-15 --- resolución y
recuento}\label{x-15--resoluciuxf3n-y-recuento}

\begin{itemize}
\tightlist
\item
  \textbf{Prosa r1: 11 «matching(s)»} (l.526, 578, 746 ×2, 1892, 1917,
  1959, 2177, 2182 ×2, 2184) \textbf{y 1 «packing»} (l.688), es decir,
  12 términos ingleses.
\item
  \textbf{Prosa r2: 0.} Las 4 apariciones que quedan están en la
  bibliografía: {[}6{]} «\ldots packings in dense graphs», {[}7{]}
  «\ldots packing of families\ldots», {[}9{]} «\ldots spread matching
  theorems», {[}12{]} «\ldots triangle packing formalization». Son
  títulos citados y deben quedar sin traducir; no hay ninguna en
  identificadores.
\item
  \textbf{Discrepancia histórica.} run\_v1.21\_r1 habló de «matching(s)
  ×12» porque contó la referencia {[}9{]} y su escaneo no detectaba
  «matchings». run\_v1.22\_r1 describió «12 matching(s) + 1 packing»
  cuando su escaneo listaba 12 apariciones en total. El recuento
  correcto es \textbf{11 + 1}. El editor tiene razón, y la corrección
  queda registrada contra el auditor (CORRECTIONS C3-01).
\item
  Las traducciones son correctas: «emparejamiento(s)» es el término
  estándar para \emph{matching} en teoría de grafos, y «empaquetamiento»
  es coherente con «empaquetamiento mixto» en el resto del texto. La
  Figura 3 sigue con «valor base», arreglado en r1.
\item
  \textbf{X-15: RESUELTO en r2.}
\end{itemize}

\subsection{\texorpdfstring{NEW-01 ---
\texttt{Model}}{NEW-01 --- Model}}\label{new-01--model}

\begin{itemize}
\tightlist
\item
  MD l.1203 «ganancia acotada de \texttt{Model}» y TeX l.1084
  \texttt{\{\textbackslash{}ttfamily\textbackslash{}small\ Model\}}. En
  el PDF, p. 35, «Model» se ve en monoespaciado (zoom a 140 dpi).
\item
  Coincide con EN y con la versión ES v1.2. Las demás partes de NEW-01
  (encabezados 5.1 y Lema 3.2, fila A.2) siguen corregidas: las vi en p.
  14, 44 y 41.
\item
  \textbf{NEW-01: RESUELTO en r2.}
\end{itemize}

\subsection{Desplazamientos de texto}\label{desplazamientos-de-texto}

\begin{itemize}
\tightlist
\item
  La expansión de los términos añade unas 2 líneas en E.1 y desplaza el
  texto en p. 55--71.
\item
  En las comparaciones r1\textbar r2 que se mostraron (p. 55, 56 y
  61--71) no hay texto perdido ni duplicado, y el corte entre páginas es
  continuo (p. 61→62, 62→63). Las fórmulas (E.1)--(F.12) y (G.1)--(G.5)
  son idénticas.
\item
  La Tabla 9 se reubica entera en p. 65 y conserva la declaración de
  X-27.
\item
  La sección de referencias empieza en p. 72 en ambas versiones, y el
  total sigue en 73 páginas.
\end{itemize}

\subsection{Texto histórico de estado (inalterado a
propósito)}\label{texto-histuxf3rico-de-estado-inalterado-a-propuxf3sito}

Las p. 1, 33, 39 y 40 (y sus equivalentes EN) dicen que las ampliaciones
de v1.22 «esperan revalidación» y presentan run\_v1.21\_r1 como la
última revalidación. Era cierto cuando se preparó r1, pero ya no
describe el estado actual (r1 PASS; r2 en curso). No es un error
matemático ni de correspondencia. Sí es un \textbf{problema editorial
real si se publica tal cual}, porque un lector entendería que C.3 no
está revisado. Lo registro como \textbf{OBSERVATION NEW-04}: debe
actualizarse antes de cualquier publicación, en ambos idiomas, lo que
rompe la identidad byte a byte del inglés y exigirá un control de
identidad nuevo. No lo corrijo (mandato §4).

\subsection{Veredicto E6 (r2)}\label{veredicto-e6-r2}

\textbf{PASS\_WITH\_OBSERVATIONS.} X-15 y NEW-01 quedan resueltos en el
artefacto; el contenido protegido es idéntico; el PDF corresponde al
TeX; las 73 páginas se revisaron. Queda una observación editorial,
NEW-04 (texto histórico de estado).

\begin{center}\rule{0.5\linewidth}{0.5pt}\end{center}

\section{NEW-03 --- Rectificación del editor
(run\_v1.22\_r2)}\label{new-03--rectificaciuxf3n-del-editor-run_v122_r2}

\subsection{Historia (se conserva, no se
reescribe)}\label{historia-se-conserva-no-se-reescribe}

\begin{itemize}
\tightlist
\item
  \texttt{01\_manuscript/v1.22\_editorial\_candidate/RESPONSE\_TO\_REVALIDATION\_v1.22.md}
  (sha256 76c86d77\ldots; el directorio r1 está intacto según E0) sigue
  diciendo dos cosas:

  \begin{itemize}
  \tightlist
  \item
    Fila X-15: «Se traducen los residuos de matching/packing en prosa
    española\ldots».
  \item
    Fila NEW-01: «Se restaura el formato del namespace Model donde
    corresponde».
  \end{itemize}
\item
  run\_v1.22\_r1 comprobó que ninguna de las dos afirmaciones era cierta
  en los artefactos españoles entregados y abrió NEW-03 (MINOR). Sus
  registros (E6\_RECORD.md, E8\_RECORD.md, FINDINGS.csv) no han cambiado
  (E0: 188/188 archivos del manifiesto).
\end{itemize}

\subsection{Rectificación en r2}\label{rectificaciuxf3n-en-r2}

En
\texttt{01\_manuscript/v1.22\_editorial\_candidate\_r2/CORRECTIONS\_AND\_HANDOFF\_v1.22\_r2.md},
§«Correcciones y rectificación», punto 3, el editor:

\begin{itemize}
\tightlist
\item
  admite que la respuesta anterior dio por terminadas las dos
  correcciones cuando no estaban aplicadas;
\item
  califica esa afirmación de incorrecta;
\item
  conserva la respuesta anterior como evidencia histórica, sin
  reescribirla.
\end{itemize}

\subsection{Evaluación}\label{evaluaciuxf3n}

\begin{itemize}
\tightlist
\item
  \textbf{Exactitud de la admisión.} Coincide con la evidencia de
  run\_v1.22\_r1: cuando se hizo la afirmación seguían sin traducir 11 +
  1 términos en la prosa, y \texttt{Model} en ES seguía en redonda.
\item
  \textbf{Conservación.} El archivo de respuesta r1 está intacto (E0).
  La rectificación se añade en un documento nuevo, sin editar el
  antiguo.
\item
  \textbf{Cierre por el artefacto.} E6\_RECORD.md verifica de forma
  independiente X-15 (prosa r2 con 0 residuos) y NEW-01 (\texttt{Model}
  monoespaciado en MD, TeX y PDF).
\item
  \textbf{Recuento.} La rectificación dice que el informe narrativo del
  auditor cuenta doce y la fuente once. Es correcto: el error de
  recuento era del auditor (CORRECTIONS C3-01), no del editor.
\end{itemize}

\textbf{NEW-03: CERRADO en r2.} Se cierra por la combinación de (a) el
artefacto corregido y verificado y (b) la rectificación explícita del
editor. La afirmación histórica y el hallazgo de r1 se conservan.

\end{document}
